\section{The main theorem in every family}\label{sec:families}

\Cref{thm:main} describes the algebraic locus of the Weil classes in one
family: it contains an explicit CM point, it is dense, it is a countable union
of closed algebraic strata, and it is the whole family as soon as it has an
interior point. The first two properties use the Weil classes; the last two
hold for every flat class on every smooth projective family
(\Cref{lem:algstrata}). This section carries the theorem to the families that
carry (F2) and
\textup{(F3$'$)}, and records what each construction of the paper
contributes there.

For (F2) the families are abelian schemes with a member of CM type. Every
Hodge class of every abelian variety lies in such a family
(\Cref{prop:cmfamily}), and (F2) is equivalent to the statement that in every
one of them the algebraic locus of a flat Hodge class that is algebraic at a
CM member has an interior point (\Cref{thm:f2propagation}). The base points
of \Cref{thm:cmbasepoint} are the CM seeds of the Weil families, and
Andr\'e's theorem makes the Weil classes of the split families seeds for every
Hodge class of CM type. For \textup{(F3$'$)} the cohomology of a variety splits
into the part reached from abelian varieties by algebraic correspondences and
a complement, and \textup{(F3$'$)} is the statement that the complement
carries no Hodge class (\Cref{prop:abpart}). On the first varieties where
\textup{(F3$'$)} is open, the squares of K3 surfaces with real
multiplication, the algebraic locus of the real multiplication contains every
CM point, is dense and stable under a group of rational Hodge isometries, and
is either everything or meagre (\Cref{thm:rmlocus}). In both cases what is
left is the interior point, and \Cref{thm:twostand} lists what the
constructions of the paper give towards it.

\subsection{Algebraic loci}\label{ssec:algloci}

\begin{definition}[Algebraic locus]\label{def:algloc}
Let $S$ be a smooth connected complex quasi-projective variety,
$f\colon\cX\to S$ a smooth projective morphism, and $\xi$ a global section of
$R^{2p}f_{*}\QQ$ over $S^{\mathrm{an}}$. The \emph{algebraic locus} of $\xi$
is
\[
  \Sigma(\xi)=\{\,s\in S:\ \xi_{s}\ \text{is the class of an algebraic cycle
  with rational coefficients on}\ \cX_{s}\,\}.
\]
\end{definition}

\begin{lemma}[The dichotomy in every family]\label{lem:algstrata}
In \Cref{def:algloc} the following hold.
\begin{enumerate}[label=\textup{(\roman*)},leftmargin=2.6em,itemsep=2pt]
\item $\Sigma(\xi)$ is the union of countably many Zariski closed subsets of
  $S$.
\item Either $\Sigma(\xi)=S$, or $\Sigma(\xi)$ is meagre in $S^{\mathrm{an}}$.
\item The following are equivalent: $\Sigma(\xi)=S$; $\Sigma(\xi)$ has an
  interior point; $\Sigma(\xi)$ is not meagre; $\xi_{s}$ is algebraic for
  every $s$ outside a countable union of proper closed analytic subsets of
  $S$. If $\Sigma(\xi)$ is dense, they are also equivalent to $\Sigma(\xi)$
  being closed.
\end{enumerate}
\end{lemma}

\begin{proof}
(i) The proof of \Cref{thm:closedlocus} applies with
$(\cA,\cS,\omega)$ replaced by $(\cX,S,\xi)$, as in \Cref{lem:spreading}:
for a tuple $\underline{P}$ of Hilbert polynomials the fibre product of the
relative Hilbert schemes $\Hilb^{P_{i}}(\cX/S)$ is projective over $S$
\cite{Gro61}, the class of a flat family of cycles is a flat section, two flat
sections over a connected base that agree at one point agree everywhere, and
the image of a proper morphism is closed. So the set of points at which
$\xi_{s}=\sum_{i}q_{i}[Z_{i}]$ with $Z_{i}$ of Hilbert polynomial $P_{i}$ is
Zariski closed, and $\Sigma(\xi)$ is the union of these strata over the
countably many pairs $(\underline{P},\underline{q})$.

(ii) and (iii). A Zariski closed subset of the irreducible variety $S$ that
has an interior point in $S^{\mathrm{an}}$ is all of $S$, and a proper one is
nowhere dense. If no stratum is $S$, then $\Sigma(\xi)$ is a countable union
of nowhere dense closed sets, hence meagre, and has no interior point, since
$S^{\mathrm{an}}$ is a Baire space. If one stratum is $S$, then
$\Sigma(\xi)=S$. This gives (ii) and the equivalence of the first three
conditions of (iii). The complement of a countable union of proper closed
analytic subsets is not meagre, by the Baire category theorem, so the fourth
condition implies the third, and it is implied by the first. Finally, a
dense set that is closed is $S$.
\end{proof}

So the dichotomy of \Cref{thm:main}(iii) and the equivalences of
\Cref{thm:main}(iv) hold for every flat class on every smooth projective
family; what \Cref{thm:main} adds for the Weil classes is a point of
$\Sigma(\xi)$ and the density, and neither is automatic, since
$\Sigma(\xi)$ may be empty.

\subsection{(F2) is propagation from CM points}\label{ssec:f2cm}

\begin{lemma}[Points of CM type are dense]\label{lem:cmdense}
Let $G$ be a connected reductive group over $\QQ$, let
$u_{0}\colon\mathrm{U}(1)\to G_{\RR}$ be a homomorphism, and let $\cD$ be
the orbit of $u_{0}$ under conjugation by $G(\RR)^{+}$, with the topology of
$G(\RR)^{+}/\mathrm{Stab}(u_{0})$. Then the points $u\in\cD$ for which
$u(\mathrm{U}(1))\subseteq T_{\RR}$ for some maximal torus $T$ of $G$ defined
over $\QQ$ are dense in $\cD$.
\end{lemma}

\begin{proof}
Let $u_{1}\in\cD$ and let $W$ be a neighbourhood of $u_{1}$. The centraliser
of the torus $u_{1}(\mathrm{U}(1))$ in $G_{\RR}$ is connected and reductive,
and a maximal torus of it defined over $\RR$ is a maximal torus $T_{1}$ of
$G_{\RR}$ containing $u_{1}(\mathrm{U}(1))$. Write $\mathfrak{g}$ for the Lie
algebra of $G$ and $\mathfrak{t}_{1}\subseteq\mathfrak{g}_{\RR}$ for that of
$T_{1}$, and choose a regular element $Y_{1}\in\mathfrak{t}_{1}$. The map
$\varphi\colon G(\RR)^{+}\times\mathfrak{t}_{1}\to\mathfrak{g}_{\RR}$,
$(g,Y)\mapsto\mathrm{Ad}(g)Y$, has differential
$(Z,Y')\mapsto[Z,Y_{1}]+Y'$ at $(1,Y_{1})$, which is onto because
$\mathfrak{g}_{\RR}=\mathfrak{t}_{1}\oplus[\mathfrak{g}_{\RR},Y_{1}]$ for a
regular semisimple $Y_{1}$. So for neighbourhoods $N_{1}$ of $1$ in
$G(\RR)^{+}$ with $N_{1}u_{1}N_{1}^{-1}\subseteq W$, and $N_{2}$ of $Y_{1}$
in $\mathfrak{t}_{1}$ consisting of regular elements, the image
$\varphi(N_{1}\times N_{2})$ contains a nonempty open subset of
$\mathfrak{g}_{\RR}$. The rational points $\mathfrak{g}(\QQ)$ are dense in
$\mathfrak{g}_{\RR}$, so that open subset contains some
$Y\in\mathfrak{g}(\QQ)$, and $Y=\mathrm{Ad}(g)Y'$ with $g\in N_{1}$ and
$Y'\in N_{2}$. As $Y$ is rational and regular semisimple, the identity
component $T$ of its centraliser in $G$ is a maximal torus defined over
$\QQ$, and $T_{\RR}=gT_{1}g^{-1}$ because the centraliser of $Y'$ has
identity component $T_{1}$. Then $u=gu_{1}g^{-1}$ lies in $W$ and
$u(\mathrm{U}(1))\subseteq T_{\RR}$.
\end{proof}

For an abelian variety the lemma is used with $G$ its Hodge group. If
$\Hg(A)$ lies in a $\QQ$-torus $T$, then $A$ is of CM type: the centraliser of
$T$ in $\End(H^{1}(A,\QQ))$ contains a commutative semisimple $\QQ$-algebra of
dimension $2\dim A$, spanned by a maximal $\QQ$-torus of $\GL(H^{1}(A,\QQ))$
containing $T$, and it lies in $\End^{0}(A)$, the commutant of $\Hg(A)$.

\begin{proposition}[A family with a CM member through every Hodge
class]\label{prop:cmfamily}
Let $A_{0}$ be a complex abelian variety and $\gamma$ a Hodge class of degree
$2p$ on $A_{0}$. There are a smooth connected complex quasi-projective variety
$S$, an abelian scheme $\pi\colon\cA\to S$, a global section $\xi$ of
$R^{2p}\pi_{*}\QQ$ that is a Hodge class at every point, and points $s_{0}$,
$c\in S$, such that $\cA_{s_{0}}\cong A_{0}$ with $\xi_{s_{0}}=\gamma$ under
this isomorphism and $\cA_{c}$ is of CM type.
\end{proposition}

\begin{proof}
Replace $\gamma$ by an integral multiple, fix a polarisation of $A_{0}$, put
$V=H_{1}(A_{0},\QQ)$, and fix a level $N\ge3$. As in \Cref{setup:shimura}, the
quotient $\cS_{g}$ of the Siegel space $\cH$ of complex structures on
$V_{\RR}$ compatible with the polarisation, by the principal congruence
subgroup of level $N$ of the stabiliser of $H_{1}(A_{0},\ZZ)$, is a smooth
quasi-projective variety carrying a universal family $f\colon\cA_{g}\to\cS_{g}$
\cite{BB66}; let $a_{0}\in\cS_{g}$ be the point of $A_{0}$ and $J_{0}\in\cH$ a
lift. Put $\mathbb{V}=R^{2p}f_{*}\ZZ$ modulo torsion, a polarised variation of
$\ZZ$-Hodge structure. By \Cref{thm:cdk} the connected component $\mathcal{C}$ of its
locus of Hodge classes that contains $(a_{0},\gamma)$ is a closed algebraic
subvariety of the total space of $\mathbb{V}$.

Let $G=\Hg(A_{0})$ and $u_{0}\colon\mathrm{U}(1)\to G_{\RR}$,
$u_{0}(e^{i\theta})=\cos\theta+J_{0}\sin\theta$, the homomorphism defining
the Hodge structure, and let
$\cD=\{gJ_{0}g^{-1}:g\in G(\RR)^{+}\}\subseteq\cH$, which is the orbit of
\Cref{lem:cmdense} when $J$ is identified with $u_{J}$. For $J=gJ_{0}g^{-1}\in\cD$ the group
$u_{J}(\mathrm{U}(1))=gu_{0}(\mathrm{U}(1))g^{-1}$ lies in $G(\RR)$, so
$\Hg(A_{J})\subseteq G$; since $\gamma$ is fixed by $G$ \cite[\S3]{Del82}, it
is a Hodge class on $A_{J}$. The map sending $J\in\cD$ to the image of $J$ in
$\cS_{g}$ together with $\gamma$ is therefore a continuous map from the
connected space $\cD$ to the locus of Hodge classes of $\mathbb{V}$, so it
takes values in $\mathcal{C}$. By \Cref{lem:cmdense} there are points $J_{k}\in\cD$
converging to $J_{0}$ at which $\Hg(A_{J_{k}})$ lies in a $\QQ$-torus, so
that $A_{J_{k}}$ is of CM type. Only finitely many irreducible components of
$\mathcal{C}$ pass through $(a_{0},\gamma)$, and every other component keeps a
positive distance from it, so one component $\mathcal{C}_{0}$ through
$(a_{0},\gamma)$ contains the image of some $J_{k}$.

Let $\rho\colon S\to\mathcal{C}_{0}$ be a resolution of singularities \cite{Hir64}:
$S$ is smooth, connected and quasi-projective, and $\rho$ is proper and
surjective. Pull $\cA_{g}$ back to $S$ along $\mathcal{C}_{0}\to\cS_{g}$. A point of
$\mathcal{C}_{0}$ is a pair $(a,v)$ with $v\in\mathbb{V}_{a}$ of Hodge type, and
$(a,v)\mapsto v$ is a flat section of the pullback of $\mathbb{V}$ to
$\mathcal{C}_{0}$; its pullback $\xi$ to $S$ is a flat section of
$R^{2p}\pi_{*}\ZZ$ modulo torsion, of Hodge type at every point. Take $s_{0}$
over $(a_{0},\gamma)$ and $c$ over the image of $J_{k}$.
\end{proof}

\begin{theorem}[(F2) is propagation from CM points]\label{thm:f2propagation}
The following are equivalent.
\begin{enumerate}[label=\textup{(\roman*)},leftmargin=2.6em,itemsep=2pt]
\item \textup{(F2)}.
\item The Hodge conjecture holds for every complex abelian variety.
\item For every abelian scheme $\pi\colon\cA\to S$ over a smooth connected
  complex quasi-projective variety and every global section $\xi$ of
  $R^{2p}\pi_{*}\QQ$, if $\xi_{s}$ is algebraic at one point $s\in S$, then
  $\Sigma(\xi)=S$.
\item The same, for the points $s$ at which $\cA_{s}$ is of CM type: if
  $\xi_{c}$ is algebraic at one point $c$ with $\cA_{c}$ of CM type, then
  $\Sigma(\xi)=S$.
\item For every $\pi$ and $\xi$ as in \textup{(iii)} with $\xi_{c}$ algebraic
  at a point $c$ with $\cA_{c}$ of CM type, $\Sigma(\xi)$ has an interior
  point.
\end{enumerate}
\end{theorem}

\begin{proof}
(i) and (ii) are equivalent by \Cref{prop:f2isab}. If (ii) holds and
$\xi_{s}$ is algebraic, then $\xi_{t}$ is a Hodge class at every $t$ by the
first part of the proof of \Cref{prop:lefvhc}, which uses only the theorem of
the fixed part and the semisimplicity of polarisable Hodge structures, and it
is algebraic by (ii); so (ii) implies (iii). Clearly (iii) implies (iv), and
(iv) and (v) are equivalent by \Cref{lem:algstrata}(iii).

Assume (iv). Let $F$ be a CM field, $n\ge2$, and consider the family of
$(F,n,\delta_{0})$-Weil type with trivial discriminant $\delta_{0}$ over its
Shimura variety with a level structure, as in \Cref{setup:cmfield} and
\Cref{setup:shimura}: an abelian scheme over a smooth connected
quasi-projective base on which the Weil classes are flat sections
(\Cref{prop:cmweil}). By \Cref{thm:main}(i) they are algebraic at a member
isogenous to a product of abelian varieties with complex multiplication by
$F$, which is of CM type. By (iv) they are algebraic at every member, so every
abelian variety of split Weil type has algebraic Weil classes, for $n\ge2$
by this argument and for $n=1$ by the theorem of Lefschetz on $(1,1)$
classes. By Andr\'e's theorem \cite[Theorem~1]{Mil20}, every Hodge class on
an abelian variety $B$ of CM type is a sum $\sum f_{\Delta}^{*}(t_{\Delta})$
of pullbacks along homomorphisms $f_{\Delta}\colon B\to A_{\Delta}$ of Weil
classes $t_{\Delta}$ on abelian varieties $A_{\Delta}$ of split Weil type; so
the Hodge conjecture holds for every abelian variety of CM type. Now let
$A_{0}$ be any abelian variety and $\gamma$ a Hodge class on it, and take
$S$, $\pi$, $\xi$, $s_{0}$, $c$ from \Cref{prop:cmfamily}. Then $\xi_{c}$ is
algebraic, so by (iv) $\xi_{s_{0}}=\gamma$ is algebraic. This is (ii).
\end{proof}

The implication from \textup{(iii)} to \textup{(ii)} is the argument of
Deligne and Andr\'e as Milne records it \cite[\S5.1, Theorem~3 and
Remark~2]{Mil20}. \Cref{thm:f2propagation} adds that only abelian schemes
and only starting points of CM type are needed, and the converse. The seeds
are explicit: for the Weil classes they are the base points of
\Cref{thm:cmbasepoint}, and for every other Hodge class they are the CM
members supplied by \Cref{prop:cmfamily}, at which the class is algebraic once
the Weil classes of the split families are. For the Weil classes of the
split family of a CM field, whose base point is of CM type, statement
\textup{(v)} is the condition of \Cref{thm:main}(iv). So the condition of
\Cref{thm:main}(iv), required in every abelian scheme with a member of CM
type at which the class is algebraic, is exactly (F2); it is not a reduction
of (F2) to something smaller, because the interior point it asks for is the
conclusion of \textup{(iii)}.

\subsection{(F3$'$) and the part reached from abelian varieties}
\label{ssec:abpart}

\begin{definition}[The abelian part]\label{def:abpart}
For a smooth complex projective variety $X$ and $p\ge0$ let
$H^{2p}_{\mathrm{ab}}(X)\subseteq H^{2p}(X,\QQ)$ be the sum of the images of
the maps
$\Gamma_{*}\colon H^{k}(B,\QQ)\to H^{2p}(X,\QQ)$,
$y\mapsto\mathrm{pr}_{X*}([\Gamma]\cup\mathrm{pr}_{B}^{*}y)$, over all
complex abelian varieties $B$, the point included, all $k$, and all algebraic
cycles $\Gamma$ on $B\times X$ with rational coefficients and of the
codimension that lands in degree $2p$.
\end{definition}

\begin{proposition}[\textup{(F3$'$)} is a statement about the complement of
the abelian part]\label{prop:abpart}
Let $X$ be a smooth complex projective variety and $p\ge0$.
\begin{enumerate}[label=\textup{(\roman*)},leftmargin=2.6em,itemsep=2pt]
\item $H^{2p}_{\mathrm{ab}}(X)$ is a sub-Hodge structure of
  $H^{2p}(X,\QQ)$, spanned by the images of finitely many of the maps
  $\Gamma_{*}$.
\item $\mathrm{Ab}^{p}(X)=\Hdg^{p}(X)\cap H^{2p}_{\mathrm{ab}}(X)$.
\item Let $T^{2p}_{\mathrm{ab}}(X)$ be a complement of
  $H^{2p}_{\mathrm{ab}}(X)$ in $H^{2p}(X,\QQ)$ as a Hodge structure. Then
  \textup{(F3$'$)} holds for $X$ in degree $2p$ if and only if
  $T^{2p}_{\mathrm{ab}}(X)$ contains no nonzero Hodge class, in particular
  whenever the Hodge group of $T^{2p}_{\mathrm{ab}}(X)$ fixes no nonzero
  vector.
\item $X$ lies in the class $\mathcal{A}$ of \Cref{prop:f3prime}(iii) if and
  only if $T^{2p}_{\mathrm{ab}}(X)=0$ for every $p$ and the odd cohomology is
  reached in the same way.
\end{enumerate}
\end{proposition}

\begin{proof}
(i) An algebraic cycle $\Gamma$ of codimension $e$ on $B\times X$, with
$\dim B=b$, induces a morphism of Hodge structures
$H^{k}(B,\QQ)\to H^{k+2e-2b}(X,\QQ)(e-b)$, so its image is a sub-Hodge
structure; a sum of sub-Hodge structures is one, and a finite-dimensional
space is spanned by finitely many of them.

(ii) Every $\Gamma_{*}\beta$ with $\beta$ a Hodge class lies in both sets, so
$\mathrm{Ab}^{p}(X)$ is contained in the intersection. Conversely take
$(B_{j},\Gamma_{j},k_{j})$, $1\le j\le r$, as in (i), and let
\[
  \Phi=\sum_{j}\Gamma_{j*}\colon
  \bigoplus_{j}H^{k_{j}}(B_{j},\QQ)(t_{j})\longrightarrow
  H^{2p}_{\mathrm{ab}}(X),
\]
with the Tate twists $t_{j}$ that make it a morphism of pure Hodge structures
of weight $2p$; it is surjective. Polarisable Hodge structures form a
semisimple category, so $\Phi$ has a section $\sigma$ that is a morphism of
Hodge structures. For a Hodge class $\gamma\in H^{2p}_{\mathrm{ab}}(X)$ the
class $\sigma(\gamma)$ is of type $(p,p)$, so each component
$\beta_{j}\in H^{k_{j}}(B_{j},\QQ)$ is a Hodge class, $k_{j}$ being even, and
$\gamma=\sum_{j}\Gamma_{j*}\beta_{j}\in\mathrm{Ab}^{p}(X)$.

(iii) A complement exists by semisimplicity, and the Hodge classes of
$H^{2p}(X,\QQ)$ are the sums of those of the two summands; by (ii) the
Hodge classes of the first summand are $\mathrm{Ab}^{p}(X)$. So
$\Hdg^{p}(X)=\mathrm{Ab}^{p}(X)$ exactly when the second has none, and a
Hodge class of $T^{2p}_{\mathrm{ab}}(X)$ is a vector fixed by its Hodge group
\cite[\S3]{Del82}. (iv) is the definition of $\mathcal{A}$ read degree by
degree.
\end{proof}

So \textup{(F3$'$)} asks nothing of the abelian part, where by (ii) and
\Cref{prop:f2isab} the Hodge conjecture reduces to (F2), and it asks of the
complement that it carry no Hodge class. A Hodge class there can be removed
only by enlarging the abelian part, that is by constructing a new algebraic
cycle on some $B\times X$, which is the step \Cref{rem:f3primeopen} isolates.
\Cref{prop:f3prime}(iii), \Cref{prop:aclosure} and \Cref{thm:simplextype}
are cases in which the complement is zero, and \Cref{prop:k3powers} and
\Cref{prop:k3ntype} cases in which the conjecture itself makes its Hodge
classes algebraic.

\subsection{The first open case of (F3$'$) in families}\label{ssec:rmlocus}

The frontier of \textup{(F3$'$)} in \Cref{rem:f3primefrontier} begins with
$S\times S$ and $S^{[2]}$ for a K3 surface $S$ with real multiplication. These
surfaces come in families over Hermitian symmetric domains, and the argument
of \Cref{thm:main} applies to them with rational Hodge isometries in place of
isogenies.

\begin{setupenv}[A family with real multiplication]\label{setup:rm}
Let $E$ be a totally real field of degree $r\ge2$ and $T$ a finite-dimensional
$\QQ$-vector space with an action of $E$ and a nondegenerate symmetric form
$q$ for which $E$ acts by self-adjoint maps, so that $q=\Tr_{E/\QQ}q_{E}$ for
a unique $E$-bilinear form $q_{E}$. Put $m=\dim_{E}T\ge3$, and suppose that
$q_{E}$ has signature $(2,m-2)$ at one real place $\sigma_{1}$ of $E$ and is
negative definite at the others, and that $(T,q)$ embeds isometrically in
$\Lambda_{\QQ}$, where $\Lambda$ is the K3 lattice; write
$N=T^{\perp}\cap\Lambda$. Let $T_{1}=T\otimes_{E,\sigma_{1}}\RR$ and
\[
  \cD_{E}=\{\,[\omega]\in\PP(T_{1}\otimes_{\RR}\CC):\ q(\omega,\omega)=0,\
  q(\omega,\bar\omega)>0\,\}^{+},
\]
one of its two components, a bounded symmetric domain of dimension $m-2$. For
$\omega\in\cD_{E}$ let $S_{\omega}$ be a K3 surface with a marking
$H^{2}(S_{\omega},\ZZ)\cong\Lambda$ carrying $H^{2,0}(S_{\omega})$ to
$\CC\omega$, which exists by the surjectivity of the period map
\cite[Chs.~6--7]{Huy16}; it is projective, since $N$ has signature
$(1,21-rm)$ and so contains a class of positive square. Each $e\in E$ acts on
$T_{1}\otimes\CC$ by the scalar $\sigma_{1}(e)$ and is self-adjoint, so it
preserves the Hodge decomposition of $T$, and $T\subseteq H^{2}(S_{\omega},\QQ)$
is a sub-Hodge structure containing $T(S_{\omega})_{\QQ}$. Write
$\tilde e_{\omega}\in H^{4}(S_{\omega}\times S_{\omega},\QQ)$ for the class of
the endomorphism $e\oplus0$ of $H^{2}(S_{\omega},\QQ)=T\oplus N_{\QQ}$, a
Hodge class, and let
\[
  \Sigma_{E}=\{\,\omega\in\cD_{E}:\ \tilde e_{\omega}\ \text{is algebraic
  for every}\ e\in E\,\}
\]
be the algebraic locus of the real multiplication. Let
$G=\Res_{E/\QQ}\SO(T,q_{E})$, so that
$G(\RR)=\SO(T_{1})\times\prod_{\sigma\ne\sigma_{1}}\SO(T_{\sigma})$ and
$G(\RR)^{+}$ acts transitively on $\cD_{E}$.
\end{setupenv}

The surfaces with real multiplication of \cite{vG08,vGS25} lie in families
of this kind, and so do the surfaces $S_{\lambda}$ of \Cref{thm:mumfordks},
for which $E$ is cubic, $m=3$ and $\cD_{E}$ is a disc.

\begin{theorem}[The algebraic locus of a real multiplication]
\label{thm:rmlocus}
Keep \Cref{setup:rm}.
\begin{enumerate}[label=\textup{(\roman*)},leftmargin=2.6em,itemsep=2pt]
\item The points $\omega$ at which the Hodge group of $T$ is commutative are
  dense in $\cD_{E}$, and they lie in $\Sigma_{E}$.
\item $\Sigma_{E}$ is stable under $G(\QQ)\cap G(\RR)^{+}$.
\item $\Sigma_{E}$ is dense in $\cD_{E}$. After a finite cover with level
  structure it is the preimage of a countable union of Zariski closed subsets
  of an arithmetic quotient of $\cD_{E}$, and either $\Sigma_{E}=\cD_{E}$ or
  $\Sigma_{E}$ is meagre.
\item The following are equivalent: $\Sigma_{E}$ is closed; $\Sigma_{E}$ has
  an interior point; $\Sigma_{E}=\cD_{E}$; the Hodge conjecture holds for
  $S_{\omega}\times S_{\omega}$ for every $\omega$ outside a countable union
  of proper closed analytic subsets of $\cD_{E}$. When they hold, the
  Hodge conjecture and \textup{(F3$'$)} hold for $S_{\omega}\times S_{\omega}$
  and $S_{\omega}^{[2]}$ at every $\omega$ at which
  $\End_{\mathrm{Hdg}}(T(S_{\omega})_{\QQ})$ is $E$ or a CM field.
\end{enumerate}
\end{theorem}

\begin{proof}
(i) The Hodge structure of $S_{\omega}$ on $T$ is given by
$u_{\omega}\colon\mathrm{U}(1)\to\SO(T_{\RR})$, acting by $z^{2}$ on
$\CC\omega$, by $z^{-2}$ on $\CC\bar\omega$ and trivially on their orthogonal
complement; it commutes with $E$ and lies in $G(\RR)$, and
$u_{g\omega}=gu_{\omega}g^{-1}$ for $g\in G(\RR)^{+}$. By \Cref{lem:cmdense}
the points at which $u_{\omega}(\mathrm{U}(1))$ lies in a maximal
$\QQ$-torus of $G$ are dense, and at them the Hodge group of $T$, the
smallest $\QQ$-subgroup whose real points contain $u_{\omega}(\mathrm{U}(1))$,
is commutative. At such a point the Hodge group $H$ of
$P=T(S_{\omega})_{\QQ}$ is commutative too, being a quotient of that of $T$.
As $P$ is an irreducible Hodge structure it is an irreducible representation
of the torus $H$, whose rational points are Zariski dense; so the commutant
$E'=\End_{\mathrm{Hdg}}(P)$ is a field with $\dim_{E'}P=1$. By \cite{Zar83}
$E'$ is totally real or CM. If it were totally real, the embedding $\tau$
through which it acts on the line $H^{2,0}(S_{\omega})$ would be real, and
that line would lie in $P\otimes_{E',\tau}\CC$, which is stable under complex
conjugation and so also contains $H^{0,2}(S_{\omega})$, giving
$\dim_{E'}P\ge2$. So $E'$ is a CM
field, the Hodge conjecture holds for $S_{\omega}\times S_{\omega}$ by
\Cref{prop:k3powers}(i), and the Hodge classes $\tilde e_{\omega}$ are
algebraic.

(ii) Let $g\in G(\QQ)\cap G(\RR)^{+}$ and $\omega\in\Sigma_{E}$. The map
$\phi=\mathrm{id}_{N}\oplus g$ is a rational isometry of
$\Lambda_{\QQ}=N_{\QQ}\oplus T$ carrying $\CC\omega$ to $\CC g\omega$, so
through the markings it is a rational Hodge isometry
$H^{2}(S_{\omega},\QQ)\to H^{2}(S_{g\omega},\QQ)$, and it is induced by an
algebraic class $Z\in H^{4}(S_{\omega}\times S_{g\omega},\QQ)$ by Buskin's
theorem \cite{Bus19}. The transpose of $Z$ induces the adjoint of $\phi$,
which is $\phi^{-1}$. As $g$ is $E$-linear,
$\phi\circ(e\oplus0)\circ\phi^{-1}=e\oplus0$, so
$\tilde e_{g\omega}=Z\circ\tilde e_{\omega}\circ Z^{t}$ is a composite of
algebraic correspondences and is algebraic.

(iii) Density follows from (i). After a finite cover with level structure the
arithmetic quotient of $\cD_{E}$ carries a family of K3 surfaces, as in
\Cref{rem:mumfordpowersopen}, and the classes $\tilde e$ are flat sections
over the family of their squares; so \Cref{lem:algstrata}(i) and (ii) apply.

(iv) The first three conditions are equivalent by \Cref{lem:algstrata}(iii)
and the density, and they follow from the fourth because $\tilde e_{\omega}$
is a Hodge class. If $\Sigma_{E}=\cD_{E}$ and
$\End_{\mathrm{Hdg}}(T(S_{\omega})_{\QQ})=E$, then each $e\in E$ acts on
$T(S_{\omega})_{\QQ}$ through the algebraic class
$\pi_{T}\circ\tilde e_{\omega}\circ\pi_{T}$, where $\pi_{T}$ is the algebraic
projector onto $T(S_{\omega})_{\QQ}$ of the proof of \Cref{prop:k3powers};
so \Cref{prop:k3powers}(iii) gives the Hodge conjecture for
$S_{\omega}\times S_{\omega}$ and $S_{\omega}^{[2]}$, and with it
\textup{(F3$'$)}. If the field is CM, \Cref{prop:k3powers}(i) gives the
same. Finally, at a very general point $T(S_{\omega})_{\QQ}=T$ and the
field is $E$, so the fourth condition follows from the third. Indeed, a
nonzero $t\in T$ of type $(1,1)$ at every point would be orthogonal to
$T_{1}\otimes\CC$, which the lines $\CC\omega$ span, and then its conjugates
under $\mathrm{Aut}(\CC)$, which permutes the summands
$T\otimes_{E,\sigma}\CC$ of $T\otimes\CC$, show $t=0$. An endomorphism of $T$
that is a Hodge class at every point commutes with every $u_{\omega}$, hence
with the subgroup they generate, a connected normal subgroup of
$\SO(T_{1})^{+}$ that is not central in any simple factor, so all of
$\SO(T_{1})^{+}$; as $T_{1}$ is absolutely irreducible under it for $m\ge3$,
the endomorphism acts on $T_{1}\otimes\CC$ by a scalar, by conjugation on
every $T\otimes_{E,\sigma}\CC$ by a scalar, and so lies in $E$. The loci
where one of the countably many other rational vectors or endomorphisms is
of Hodge type are therefore proper closed analytic subsets.
\end{proof}

On the families of van Geemen and Sch\"utt with real multiplication by
$\QQ(\zeta_{7}+\zeta_{7}^{-1})$ and $\QQ(\zeta_{9}+\zeta_{9}^{-1})$ the
real multiplication is induced by explicit cycles \cite[Theorem~1.2 and
Section~4.8]{vGS25}, and on the double planes branched along six lines that
have real multiplication, necessarily by a quadratic field, by Schlickewei's
cycles \cite{Sch10}. Wherever such members fill an open subset of
$\cD_{E}$, (iv) extends the real multiplication to every member. For the family of the surfaces $S_{\lambda}$,
\Cref{thm:rmlocus} gives the real multiplication at every CM point and on a
dense set stable under $G(\QQ)$, and \Cref{thm:mumfordks}(ii) gives it at
$S_{\lambda}$ from the Kuga--Satake class; whether $\Sigma_{E}$ is closed
there is open. The Hodge isometries of (ii) move the real multiplication
from one surface to another but never create it on a surface where it is
not yet algebraic, for the reason given in \Cref{rem:mumfordks}: the
composites that return to $T(S_{\lambda})$ span only the multiquadratic
subfield of $E$, which is $\QQ$ for a cubic field.

\subsection{What the constructions reach}\label{ssec:reach}

\begin{theorem}[Where (F2) and (F3$'$) stand]\label{thm:twostand}
\leavevmode
\begin{enumerate}[label=\textup{(\roman*)},leftmargin=2.6em,itemsep=3pt]
\item \emph{What carries over.} In every abelian scheme of
  \Cref{thm:f2propagation} and every family of \Cref{setup:rm}, the algebraic
  locus is a countable union of closed algebraic strata and is either
  everything or meagre (\Cref{lem:algstrata}, \Cref{thm:rmlocus}); in the
  Weil families and the families of \Cref{setup:rm} it is dense and contains
  CM points, explicit ones in the Weil families (\Cref{thm:main}, \Cref{thm:rmlocus}); every Hodge class
  of every abelian variety lies in a family with a CM member
  (\Cref{prop:cmfamily}).
\item \emph{(F2).} It holds if and only if the algebraic locus of every flat
  Hodge class algebraic at a CM member of an abelian scheme has an interior
  point (\Cref{thm:f2propagation}). The constructions that supply an
  interior point are: an object with injective semiregularity map, or
  satisfying the flatness condition, at one point
  (\Cref{thm:semiregclosure}, \Cref{thm:cmpropagation}, \Cref{prop:p2flat});
  a bound on the Hilbert data of representatives on a Zariski dense set
  (\Cref{thm:degreebound}, \Cref{thm:modpdensity}), which is equivalent to
  its conclusion (\Cref{thm:p1equivalent}); and the Lefschetz standard
  conjecture for the total space over a curve through a base point
  (\Cref{cor:weilequiv}, \Cref{cor:mumfordequiv}), which is equivalent to it.
  Markman's sheaves are objects of the first kind, and with them he proves
  the algebraicity of the Weil classes on every abelian fourfold of Weil
  type and on the split sixfolds of every imaginary quadratic field
  \cite{Mar25a,Mar25b}; with \cite{MZ99} this gives the Hodge conjecture for
  abelian varieties of dimension at most five. The open cases named in
  \Cref{sec:scope} include the non-split Weil sixfolds of imaginary
  quadratic fields, the Weil eightfolds of quartic CM fields, where an
  object of the first kind needs
  $\dim\Ext^{2}(E,E)\ge88$ (\Cref{cor:quarticleast}), and the squares of
  Mumford fourfolds, where the classes are algebraic at every CM point
  (\Cref{thm:mumfordcm}) and on every member once one Kuga--Satake class is
  algebraic (\Cref{thm:mumfordks}), or once one complex at one CM point has
  $\Ext^{2}$ of dimension $119$ (\Cref{thm:mumfordnumerical}).
\item \emph{(F3$'$).} It holds if and only if, for every variety and every
  degree, the complement of the abelian part carries no Hodge class
  (\Cref{prop:abpart}). The complement is zero on the class $\mathcal{A}$,
  which contains the Fermat and Delsarte hypersurfaces of every dimension and
  the cyclic covers branched along them (\Cref{prop:aclosure},
  \Cref{cor:delsarteall}), and it carries no Hodge class on the varieties of
  \Cref{prop:f3primesmall}, \Cref{prop:f3primefibration} and
  \Cref{prop:f3primedominant}, where \textup{(F3$'$)} is proved, and of
  \Cref{prop:k3powers} and \Cref{prop:k3ntype}, where the conjecture itself
  is.
  On the first open varieties, $S\times S$ and $S^{[2]}$ for a K3 surface
  with real multiplication, the algebraic locus of the real multiplication
  is dense, contains every CM point and is either the whole family or meagre
  (\Cref{thm:rmlocus}); the real multiplication is algebraic on the members
  of the families of \cite{vGS25} and on the quadratic double planes of
  \cite{Sch10}, and whether the locus is closed is open for the family of
  $S_{\lambda}$. Beyond the varieties of abelian type, a Hodge
  class off the abelian part is removed only by a new algebraic cycle on a
  product $B\times X$ (\Cref{rem:f3primeopen}).
\end{enumerate}
Neither (F2) nor \textup{(F3$'$)} is proved here, and by \Cref{thm:final}
together they are the Hodge conjecture.
\end{theorem}

\begin{proof}
Each statement is the result cited with it. For (ii), Markman's theorems are
recalled in the introduction and their reach in \Cref{sec:scope}, where the
non-split sixfolds and the fields of degree at least four are among the open
entries, and \Cref{thm:quarticrank} and
\Cref{cor:quarticleast} give the number $88$. For (iii), the known cases of
the real multiplication are those recalled in \Cref{rem:mumfordks}.
\end{proof}
